\RequirePackage{pdfmanagement}
\SetKeys[document/metadata]{pdfversion=1.7, pdfstandard=A-2u, lang=en-GB}
\input{preamble/pdfA-compliance/glyphtounicode}
\documentclass[11pt,a4paper]{report}

\newcommand*{\PaperTitle}{The Hopf branches \(j\ge4\) of the Becker--Döring equations with atomization}
\newcommand*{\PaperTitlePlain}{The Hopf branches j >= 4 of the Becker-Doring equations with atomization}
\newcommand*{\PaperShortTitle}{Higher Hopf branches in Becker--Döring with atomization}
\newcommand*{\PaperSubtitle}{chemistry-open-problems \textperiodcentered{} 25 October 2026 \textperiodcentered{} partial result, not yet verified}
\newcommand*{\PaperCredits}{Solver: Claude Opus 5.5 (effort xhigh); editing: GPT-6-Astra (effort xhigh); verification pending.}
\newcommand*{\PaperKeywords}{Becker--Döring equations, nucleation, Hopf bifurcation, first Lyapunov coefficient, interval arithmetic}
\input{chemop-preamble}
\addbibresource{references.bib}

\begin{document}
\makepapertitle

\begin{paperabstract}
Pego and Velázquez proved that the Becker--Döring equations with atomization undergo Hopf bifurcations on branches
\(\kappa=\kappa_j(M)\), \(j=1,2,\dots\), for large maximal cluster size \(M\). They left open whether the bifurcating
solutions are stable. Earlier work in this repository addressed this for \(j\le3\). Here we treat \(j\ge4\).
(i) In the limit of infinite atomization rate the linearization becomes an explicit matrix whose unstable eigenvalue
pairs we count rigorously for every \(M\le1150\). The \(j\)-th pair appears at a real value \(M^*_j\in(M_j-1,M_j)\), with
\(M_j=142,210,289,379,480,591,713,844,985,1135\) for \(j=4,\dots,13\). In a chart that is regular at \(K=\infty\) we
follow the \(j\)-th Hopf curve from that point to the regime of Pego and Velázquez. For \(j=4,\dots,13\) the whole
curve is certified, so branch \(j\) has exactly one Hopf point for every integer \(M\ge M_j\) and none below.
(ii) For \(j=4,\dots,13\) and every \(M\ge M_j\) we certify, in ball arithmetic: simplicity and transversality, the
absence of other eigenvalues on the imaginary axis, the presence of \(j-1\) unstable pairs (counted as zeros of the
characteristic function), and the sign of the first Lyapunov coefficient \(\ell_1\). The coefficient is negative,
except at the two onsets \((j,M)=(4,142)\) and \((13,1135)\), where \(\ell_1>0\). The bifurcation is therefore
supercritical on the centre manifold, apart from these two onset points, and the bifurcating orbits are unstable in
the full system.
(iii) For every \(j\) we prove that \(\lim_{M\to\infty}\ell_1/M\) equals the negative closed-form value predicted
formally in the earlier work. Hence every branch is supercritical for all sufficiently large \(M\), with a threshold
that is not explicit.
\end{paperabstract}

\section*{The question}

R. L. Pego and J. J. L. Velázquez~\autocite{PV2020} study the Becker--Döring equations with atomization, sizes
\(1\le\ell\le M\), atomization rate \(K\). Their Theorem~3.3 gives a Hopf bifurcation at \(\kappa=\kappa_j(M)\),
\(\kappa=K\sqrt M\), for \(j\le k\) and \(M>M_{0,k}\) (not explicit). Right after it (§3.3) they write:
\begin{openquote}
``We have not managed to determine analytically whether the bifurcating solutions are stable (the supercritical case)
or not.''
\end{openquote}
Result~8 of this repository~\autocite{Result8} and its extension to every \(M\)~\autocite{AllM} settled branches
\(j=1,2,3\). Branch 1 is supercritical with orbitally stable orbits for every \(M\ge16\) and subcritical at \(M=15\).
Branches 2 and 3 are supercritical on the centre manifold, with unstable orbits. This paper treats \(j\ge4\).

\section{Setting and statements}

We use the notation of~\autocite{Result8,AllM}. The unknowns are \(n_\ell\), the constant equilibrium is \(n\equiv
A=1+K\), and \(B=B(K,M)\) is the Jacobian (PV (27)). The dynamics is exactly \(\dot v=Bv+v_1Sv\), restricted to the
mass hyperplane \(H\). \(\ell_1\) is the first Lyapunov coefficient of the restricted system (Kuznetsov's projection
formula~\autocite{Kuznetsov2004}). Further notation:
\begin{itemize}
\item \(\varepsilon=M^{-1/2}\), \(z=M(\varphi-1)\), \(\lambda=(A-\varphi^{-1})(\varphi-1)\);
\item \(F(\varphi)\) is the characteristic function of PV (71)--(76); by PV Prop.~5.3 its zeros correspond to the
eigenvalues of \(B\);
\item \(t_j\) is the \(j\)-th positive root of \(\tan t=t\) with \(\cos t<0\), and \(\kappa_j^0=(\sqrt{1+t_j^2}-1)^{1/2}\)
(PV (89)--(91)).
\end{itemize}
For real \(M\) put \(\varphi^M=\exp(M\log\varphi)\). The Hopf set is \(\mathcal H=\{(M,K,\varphi):M>3,\ K>0,\
\operatorname{Im}\varphi>0,\ F=0,\ \operatorname{Re}\lambda=0\}\). For integer \(M\) its points are exactly the pairs
\(\pm i\omega\) of eigenvalues of \(B\).

\begin{definition}
Branch \(j\), \(\Gamma_j\), is the connected component of \(\mathcal H\) that contains Pego and Velázquez's
\(j\)-th Hopf curve near \(M=\infty\). That curve is obtained from the limit root \((\kappa_j^0,it_j)\) by the implicit
function theorem (PV Lemma~7.3, Prop.~9.5). Its point at an integer \(M\) is written \(\kappa_j(M)\).
\end{definition}

\begin{theorem}[note={branches \(4\le j\le13\)}]\label{thm:A}
Let \(4\le j\le13\) and \(M_j\) as in \zcref{tab:results}.
\begin{enumerate}
\item \(\Gamma_j\) is a \(C^1\) arc. It is the graph of a function \(M\mapsto(K,\varphi)\) on \((M^*_j,\infty)\), with
\(M_j-1<M^*_j<M_j\) and \(K\to\infty\) as \(M\downarrow M^*_j\). So branch \(j\) has exactly one Hopf point at every
integer \(M\ge M_j\) and none at \(M<M_j\).
\item At every \(\kappa_j(M)\), \(M\ge M_j\):
  \begin{itemize}
  \item \(\pm i\omega\) are simple eigenvalues of \(B\) and \(\operatorname{Re}d\lambda/d\kappa>0\);
  \item \(0,\pm i\omega\) are the only eigenvalues on \(i\mathbb R\);
  \item exactly \(2j-2\) zeros of \(F\) with \(|\varphi|\ge1\), counted with multiplicity, give eigenvalues with
  \(\operatorname{Re}\lambda>0\).
  \end{itemize}
\item \(\ell_1<0\) at every \(\kappa_j(M)\), \(M\ge M_j\), except at \((j,M)=(4,142)\) and \((13,1135)\), where
\(\ell_1=+5.896\ldots\times10^{-15}\) and \(\ell_1=+2.316\ldots\times10^{-20}\) (normalization \(q^{\mathsf H}q=1\)).
\end{enumerate}
\end{theorem}

\begin{corollary}
For \(4\le j\le13\) and every integer \(M\ge M_j\), the Hopf bifurcation at \(\kappa_j(M)\) of the system on \(H\) is
supercritical on its two-dimensional centre manifold, except at \((j,M)=(4,142)\) and \((13,1135)\), where it is
subcritical. In the full
hyperplane \(H\) the bifurcating periodic orbits are unstable.
\end{corollary}

\begin{theorem}[note={\(K\to\infty\)}]\label{thm:B}
Let \(C(M)=\lim_{K\to\infty}B/A\) and let \(u(M)\) be half the number of its eigenvalues in \(\operatorname{Re}\nu>0\),
counted with multiplicity. For every integer \(3\le M\le1150\):
\begin{itemize}
\item \(C(M)\) has no eigenvalue on \(i\mathbb R\setminus\{0\}\), and \(0\) is simple;
\item \(u(M)=\#\{j:M_j\le M\}\), with the onset sequence
\[
(M_j)_{j\le13}=(15,44,87,142,210,289,379,480,591,713,844,985,1135);
\]
\item for all large \(K\), \(B(K,M)|_H\) has exactly \(u(M)\) pairs of non-real eigenvalues with
\(\operatorname{Re}\lambda>0\).
\end{itemize}
The \(j\)-th pair of \(C\) crosses \(i\mathbb R\) from left to right at \(M=M^*_j\in(M_j-1,M_j)\), \(j\le13\).
\end{theorem}

\begin{theorem}[note={every branch, large \(M\)}]\label{thm:C}
For every \(j\ge1\), along \(\Gamma_j\),
\[\lim_{M\to\infty}\frac{\ell_1(q_1=1)}M=L(t_j)<0,\]
where, with \(R=\sqrt{1+t^2}\),
\[L(t)=-\frac{R(R-1)(7R^2-34R+23)}{t(R+1)(5R-3)\big((R-1)^3+R+1\big)}.\]
Hence there is \(\widetilde M_j\) with \(\ell_1<0\) at \(\kappa_j(M)\) for every integer \(M\ge\widetilde M_j\). With PV
Theorem~3.2, the bifurcation on every branch is supercritical on the centre manifold for all large \(M\). The orbits
are stable for \(j=1\) and unstable for \(j\ge2\).
\end{theorem}

The closed form \(L\) was obtained in result~8 by a formal multiple-scales reduction to a renewal equation. All-M
confirmed it to 53 digits for \(j\le3\). \zcref{thm:C} turns it into a theorem for every \(j\), given all-M's
\(\varepsilon\)-regular formula for \(\ell_1\).

\section{The limit \texorpdfstring{\(K\to\infty\)}{K to infinity}}

\begin{lemma}\label{lem:split}
\(B=AC+E\) with \(C,E\) independent of \(K\), where
\begin{align*}
(Cv)_1&=-(M+2)v_1-\sum_{\ell=2}^{N}v_\ell+Mv_M,\\
(Cv)_\ell&=v_{\ell-1}-v_\ell\quad(2\le\ell\le N),\\
(Cv)_M&=v_1+v_N-v_M.
\end{align*}
\end{lemma}
\begin{proof}
Insert \(K=A-1\) in PV (27): \(1-K=-A+2\), \(-K=-A+1\), \(MK+1=AM+1-M\). The remainder \(E\) does not depend on \(K\).
\end{proof}

\begin{lemma}\label{lem:det}
Let \(G(\varphi)=-\varphi^{M-2}(\varphi^3+M\varphi^2-2M\varphi+M)+(M+1)\varphi-M\). Then
\(\det(\nu I-C)=-G(1+\nu)/\nu\). Moreover \(G(1)=G'(1)=0\), and \(R=G/(\varphi-1)^2\) has integer coefficients and
degree \(M-1\).
\end{lemma}
\begin{proof}
Let \(\varphi=1+\nu\ne0\). Put \(w_\ell=\varphi^{1-\ell}\) for \(\ell\le N\) and \(w_M=(1+\varphi^{2-M})/\varphi\). Then
rows \(2,\dots,M\) of \((C-\nu)w\) vanish, and row 1 equals some scalar \(g(\nu)\). Replace the first column of the
identity by \(w\); this matrix has determinant \(1\). The first column of \((C-\nu)W\) is \(g\,e_1\). The minor of
\(C-\nu\) without row and column 1 is lower triangular with diagonal \(-\varphi\). Hence
\(\det(C-\nu)=(-\varphi)^{M-1}g\), and a geometric sum gives \(\varphi^{M-1}(\varphi-1)g=G(\varphi)\).
\end{proof}

\begin{lemma}\label{lem:roots}
\(G\) has no root with \(\operatorname{Re}\varphi\ge1\) and \(|\varphi|\ge3\), and \(G<0\) on \((1,\infty)\).
\end{lemma}
\begin{proof}
Let \(r=|\varphi|\ge3\) and \(\operatorname{Re}\varphi\ge1\). Then \(|\varphi+M|\ge M+1\), so
\(|\varphi^3+M\varphi^2-2M\varphi+M|\ge(M+1)r^2-2Mr-M\ge r^2(2M/9+1)\). Hence
\(|G|\ge r^M(2M/9+1)-(2M+1)r>0\).

On \((1,\infty)\) we have \(-G=h\) with \(h(1)=h'(1)=0\) and \(h''>0\), so \(h>0\) there.
\end{proof}

By \zcref{lem:det,lem:roots} the eigenvalues of \(C\) are \(0\) and \(\varphi-1\), for the roots \(\varphi\) of \(R\).
Those with \(\operatorname{Re}\nu\ge0\) lie in the half-disc \(\{\operatorname{Re}\varphi\ge1,|\varphi-1|<4\}\). The
unstable ones are non-real. Since \(B/A=C+E/A\) depends continuously on \(1/A\), and the zero eigenvalue of \(B\) is
simple for every \(K\) (PV Lemma~3.1), the count of \(C\) is also the count of \(B\) for large \(K\).

\paragraph{Certificate (\texttt{kinf.py}).}
\begin{itemize}
\item \emph{K2.} For every \(3\le M\le1150\), the winding number of \(G\) is computed on the boundary of the half-disc,
indented around the double zero \(\varphi=1\). Each edge is subdivided until a mean-value enclosure of \(G\) excludes
\(0\) (all-M Lemma~4.1).
\item \emph{K3.} The resulting \(u(M)\) is nondecreasing, with unit jumps exactly at the \(M_j\) of \zcref{thm:B}.
\item \emph{K4.} A Krawczyk test~\autocite{Krawczyk1969,Moore1977} encloses the zero \((M^*_j,\eta_j)\) of
\(G(1+i\eta;M)\) for real \(M\), and certifies \(d\operatorname{Re}\varphi/dM>0\) there.
\end{itemize}
This proves \zcref{thm:B}.

\section{A chart regular at \texorpdfstring{\(K=\infty\)}{K = infinity}}

Put \(s=1/A\). Dividing PV (71)--(75) by \(A^3\) gives exactly (\texttt{check\_onset\_formula.py}, symbolic):
\[F/A^3=f_1\big(H_1-s^{M-1}\varphi^{-M}H_2\big)-\big(G_1-\varphi^MG_2\big)f_2,\]
where
\begin{gather*}
f_1=\varphi^M+s(\varphi-1),\qquad f_2=s-\varphi+s^{M-1}\varphi^{1-M},\\
G_1=\varphi(s-\varphi)\big((1-s)\varphi+(M(1-s)+s)(\varphi-1)\big),\\
G_2=(s-\varphi)\big(M(\varphi-1)+(1-s)\varphi+\varphi(\varphi-1)\big),\\
H_1=(\varphi-1)\big(s(1-s)+(M(1-s)+s)(s-\varphi)\big),\\
H_2=(\varphi-1)\big((s+M\varphi)(s-\varphi)+s(1-s)\varphi\big).
\end{gather*}
Also \(\lambda/A=(1-s/\varphi)(\varphi-1)\). The system \(\operatorname{Re}F/A^3=\operatorname{Im}F/A^3=
\operatorname{Re}\lambda/A=0\) in \((s,M,\varphi)\) is \(C^1\) up to \(s=0\). At \(s=0\) it is \(-\varphi^3G(\varphi)=0\),
\(\operatorname{Re}\varphi=1\), so its zero there is \((M^*_j,1+i\eta_j)\).

\texttt{onset\_sweep.py} certifies the \(j\)-th Hopf curve from \(s=0\) upwards in two pieces. On
\(s\in[0,2\cdot10^{-4}]\) it is a graph over \(s\) (parametric Krawczyk boxes, with \(dM/ds>0\)). Beyond that it is a
graph over real \(M\), up to past the start of the \(\varepsilon\)-chart. Each box carries a uniqueness statement.
Two junction checks place a sub-box inside the neighbouring certified box: between the two pieces, and between the
second piece and the last box of the \(\varepsilon\)-chart.

\section{Proof for branches 4--13}

\paragraph{The arc.} \(\Gamma_j\) is covered by three pieces:
\begin{itemize}
\item the \(\varepsilon\)-chart, \(\varepsilon\in[0,\varepsilon_c]\) (all-M's parametric Krawczyk boxes, which start at
PV's limit root);
\item the \(M\)-graph of the onset chart;
\item the \(s\)-graph of the onset chart.
\end{itemize}
Each piece is a graph with uniqueness in its boxes, and the junctions give common sub-arcs. The union is therefore one
\(C^1\) arc on which \(M\) is strictly monotone, with \(M\in(M^*_j,\infty)\). The arc is open in \(\mathcal H\), by
uniqueness in the boxes. It is closed in \(\mathcal H\): its closure adds only the two ends \(M=\infty\) and
\(K=\infty\). Hence it is a component of \(\mathcal H\), namely \(\Gamma_j\). This proves part 1.

\paragraph{Integers \(M\ge M_c\)} (with \(M_c^{-1/2}\le\varepsilon_c\)). On every \(\varepsilon\)-box,
\texttt{l1\_sweep.py} certifies the Hopf box, simplicity, transversality, \(\operatorname{Im}z>0\), \(f(\varphi)\ne0\),
and \(\ell_1(q_1=1)/M<0\) (all-M §5, mean value form). \texttt{verify\_sweep.py} re-proves every line.

\texttt{h2cover.py} applies all-M Theorem~4.5 on merged parameter boxes. Localization, monotonicity, the region lemma,
the level chains and the winding counts \(N(P)=2j+2\), \(N(P_-)=4\) are certified for every \((\varepsilon,\kappa)\) in
the union of the \(\varepsilon\)-ranges times the hull of the \(\kappa\)-ranges. The contour parameters are
\(R_0=2\pi j+7\) and \(Y=R_0+0.6\), with level values \(-0.1\) and \(0.25\). Near \(\varepsilon=0\), for \(j\ge5\), the
narrower values \(\mp\frac12\log(t_{j\pm1}/t_j)\pm0.02\) are needed, because the neighbouring zeros approach the axis
like \(\log(t_{j\pm1}/t_j)\).

\paragraph{Integers \(M_j\le M<M_c\).} \texttt{pointwise.py} certifies, at each \(M\):
\begin{itemize}
\item result~8's Krawczyk box for the Hopf point at integer \(M\), with simplicity, transversality and \(\omega>0\);
\item result~8's closed form of \(\ell_1\), and its sign;
\item all-M's Krawczyk box at the ball \(\varepsilon=M^{-1/2}\), which contains the first box, with
\(\operatorname{Im}z>0\) and \(f(\varphi)\ne0\);
\item all-M Theorem~4.5 at that point.
\end{itemize}
\texttt{pointwise.py --verify} re-proves every line. \texttt{link\_check.py} shows that each box lies on the certified
curve. This proves parts 2 and 3.

\paragraph{Cross-checks.} On \(M_c\le M\le3M_c\), result~8's \(\ell_1\) lies inside the \(\varepsilon\)-chart enclosures
(\texttt{check\_overlap.py}). Result~8's verifier's sparse \(O(M)\) code shares no code with ours. At \(M=142,143\)
(\(j=4\)) and \(M=1135,1136\) (\(j=13\)) it reproduces \(\ell_1\) to \(10^{-16}\), including the positive signs
(\texttt{xcheck\_sparse.py}). The parameters of every run are listed in
\texttt{certificates/logs/run\_branch\_params.txt}.

\paragraph{Instability.} At \(\kappa_j(M)\) the matrix \(B|_H\) has an eigenvalue \(\lambda_u\) with
\(\operatorname{Re}\lambda_u>0\). As \(\kappa\to\kappa_j(M)\), the monodromy matrix of the bifurcating orbit tends to
\(\exp(2\pi B|_H/\omega)\). So for \(\kappa\) close to \(\kappa_j(M)\) the orbit has a Floquet multiplier of modulus
\(>1\).

\section{The large-\texorpdfstring{\(M\)}{M} limit on every branch}

\texttt{limit\_symbolic.py} runs the algorithm of all-M's regular formula at \(\varepsilon=0\), symbolically.
\begin{itemize}
\item \emph{S1.} The result agrees with the ball-arithmetic code at \(\varepsilon=0\) for \(j\le8\).
\item \emph{S2.} With the Hopf relations \(e^{it}=-(1+it)/R\), \(R=1+\kappa^2=\sqrt{1+t^2}\), in the rational
parametrization \(R=(4\tau^4+1)/(4\tau^2)\), \(t=(4\tau^4-1)/(4\tau^2)\), \(\kappa=(2\tau^2-1)/(2\tau)\), the result is
identically equal to \(L\).
\item \emph{S3.} Every quantity the algorithm divides by is a rational function of \(\tau\) whose numerator has no root
\(\tau\ge2\). All Hopf points have \(\tau>2.12\).
\item \emph{S4.} \(L<0\) for \(\tau\ge2\).
\end{itemize}
By PV Lemmas~7.1--7.3 the limit roots are simple and cross transversally, so \(\Gamma_j\) is a \(C^1\) curve in
\(\varepsilon\in[0,\varepsilon_j)\). The regular formula is continuous there and equals \(\ell_1(q_1=1)/M\) for
\(\varepsilon=M^{-1/2}\). Its value at \(\varepsilon=0\) is \(L(t_j)\). PV Theorem~3.2 gives hypothesis H2 and \(j-1\)
unstable pairs at \(\kappa_j(M)\) for \(M>M_{0,j}\).

\section{Results}

\begin{table}[h]
\centering\small
\caption{Certified results per branch. \(M^*_j\): approximate onset at \(K=\infty\); exact enclosures are in \texttt{kinf.log}, K4. Last column: sign of
\(\ell_1\) at the first integer \(M_j\).}\label{tab:results}
\begin{tabular}{@{}rlrrlrl@{}}
\toprule
\(j\) & \(M^*_j\) & \(M_j\) & \(M_j-M^*_j\) & \(\varepsilon_c\) & integer \(M\) & \(\ell_1(M_j)\)\\
\midrule
4 & 141.98176056 & 142 & 0.01824 & 0.06610 & 142--228 & \(+\)\\
5 & 209.62894611 & 210 & 0.37105 & 0.05456 & 210--335 & \(-\)\\
6 & 288.73857833 & 289 & 0.26142 & 0.04653 & 289--461 & \(-\)\\
7 & 378.86583449 & 379 & 0.13417 & 0.04059 & 379--606 & \(-\)\\
8 & 479.65992358 & 480 & 0.34008 & 0.03609 & 480--767 & \(-\)\\
9 & 590.83427770 & 591 & 0.16572 & 0.03252 & 591--945 & \(-\)\\
10 & 712.14838889 & 713 & 0.85161 & 0.02962 & 713--1139 & \(-\)\\
11 & 843.39613168 & 844 & 0.60387 & 0.02721 & 844--1350 & \(-\)\\
12 & 984.39788852 & 985 & 0.60211 & 0.02519 & 985--1575 & \(-\)\\
13 & 1134.99502237 & 1135 & 0.00498 & 0.02346 & 1135--1816 & \(+\)\\
\bottomrule
\end{tabular}
\end{table}

\paragraph{Near the onsets} (observation). Along \(\Gamma_j\), as \(\kappa\to\infty\), the computed values follow
\(\widehat L\kappa^3\approx c_j+a_j\kappa\), with \(c_j\approx-5\) and a small \(a_j>0\). Also
\(M-M^*_j\approx C_j/\kappa\). So \(\ell_1\) changes sign once, in a window of width \(w_j\approx C_ja_j/|c_j|\) above
\(M^*_j\); we find \(w_j\approx0.03\)--\(0.05\) for \(j=4\)--\(7\). An integer onset is subcritical exactly when
\(M_j-M^*_j<w_j\). Up to \(j=13\) this happens exactly at \(j=1\) (\(0.107\)), \(j=4\) (\(0.018\)) and \(j=13\) (\(0.005\)),
in agreement with the certified signs. The next predicted case among \(j\le40\) is \(j=40\)
(\(M_{40}-M^*_{40}\approx0.0004\)), which we did not compute. Empirically \(w_j\approx0.18/j\), so \(\sum_jw_j\) diverges:
if the fractional parts of \(M^*_j\) are spread out, infinitely many subcritical onsets are to be expected. This is a
heuristic, not a claim.

\section{What is not proved}

\begin{itemize}
\item Branches \(j\ge14\) for all \(M\). For them only \zcref{thm:C} holds, with a threshold that is not explicit. No
argument uniform in \(j\) with explicit constants is given.
\item That \(M_j\) is the first \(M\) with \(j\) imaginary-axis crossings for \emph{some} \(K\): \zcref{thm:A} concerns
\(\Gamma_j\). Other components of \(\mathcal H\) are not excluded, and \zcref{thm:B} holds only for \(K\ge K_0(M)\),
with \(K_0(M)\) not explicit.
\item That the \(2j-2\) unstable zeros are simple.
\item All-M's \(\varepsilon\)-regular formula and Theorem~4.5 are used as stated there; that work is not yet verified.
\end{itemize}

\section{Computer verification}

All certificates use Arb~\autocite{Johansson2017} through python-flint 0.9, at 192--256 bits. Each script prints
PASS/FAIL per claim and exits non-zero on failure. The scripts, run times and logs are listed in \texttt{README.md}
and \texttt{certificates/logs/}.

\paperreferences
\end{document}
